% =============================================================================
%  Section 6 — Conclusion and outlook
% =============================================================================
\section{Conclusion and outlook}
\label{sec:conclusion}

This study first shows that not all rounds are equally easy to
predict. In pistol rounds, the tightly bounded economy on both sides
makes the outcome largely random (AUC~$\approx 0.56$). In post-pistol
rounds, on the contrary, the asymmetry created by the previous round
allows an almost deterministic prediction (AUC~$\approx 0.86$). The
normal regime sits in between (AUC~$\approx 0.71$). Within each
regime, the immediate state of the round — equipment, money, weapons,
utility — carries most of the predictive signal. Global team
statistics (HLTV rating, recent form) and the map effect add little
in comparison: they act as secondary levers that improve models only
marginally.

\paragraph{Outlook.}
Three natural directions extend this work:
\begin{itemize}
  \item \emph{Models}: tuning the logistic regression penalty
    ($\ell_1$, $\ell_2$, ElasticNet), non-linear models (XGBoost).
  \item \emph{Features}: enriching the dataset with individual player
    statistics, matchups, per-player form.
  \item \emph{Vision}: integrating \emph{live} signals (player
    positions, utilities thrown, strategy currently in use) for
    real-time prediction.
\end{itemize}

% =============================================================================
